% !TEX root = ../../../proposal.tex

\section{Compositional Diffusion Models}
\label{sec:compositional_diffusion}

\subsection{Product-of-Experts Composition}

Given two conditions $c_1$ and $c_2$, the na\"{i}ve approach is to condition a single
model on a joint prompt $c = (c_1, c_2)$. An alternative is to define the composed
distribution as a \emph{product of experts}~\citep{product_of_experts}:
\begin{equation}
    \hat p(x \mid c_1, c_2) \propto \frac{p(x \mid c_1)\,p(x \mid c_2)}{p(x)},
    \label{eq:poe_dist}
\end{equation}
which treats each conditional as an independent ``expert'' and combines them
multiplicatively. This factorisation is exact under the assumption that $c_1$ and $c_2$ are
conditionally independent given $x$. Taking log-gradients of Equation~\ref{eq:poe_dist}
gives the corresponding composed score:
\begin{equation}
    \nabla_x \log \hat p(x \mid c_1, c_2)
    = \nabla_x \log p(x \mid c_1)
    + \nabla_x \log p(x \mid c_2)
    - \nabla_x \log p(x).
    \label{eq:poe_score}
\end{equation}
Substituting the noise-predictor representation of the score
(Equation~\ref{eq:score_noise}) and incorporating per-concept CFG weights $w_1, w_2$,
Liu et al.~\citep{liu2022compositional} obtained the practical AND composition operator:
\begin{equation}
    \tilde\varepsilon_{\mathrm{AND}}
    = \varepsilon_\theta(x_t, t, \varnothing)
    + w_1\bigl[\varepsilon_\theta(x_t, t, c_1) - \varepsilon_\theta(x_t, t, \varnothing)\bigr]
    + w_2\bigl[\varepsilon_\theta(x_t, t, c_2) - \varepsilon_\theta(x_t, t, \varnothing)\bigr].
    \label{eq:and_operator}
\end{equation}
This requires no model retraining: three forward passes through a single pretrained
model are combined at each denoising step. Liu et al.\ also define a NOT operator by
negating the guidance term of the concept to be excluded, enabling structured logical
composition of multiple concepts. Phase~1 and Phase~3 both use Equation~\ref{eq:and_operator}
as the composition rule, with Phase~1 examining under what conditions it succeeds or fails
and Phase~3 applying it to co-morbid disease synthesis.

\subsection{When Composition Fails: Projective Mechanisms}
\label{subsec:bg_projective}

Despite its empirical success in many settings, PoE composition has a fundamental limitation.
Bradley~\citep{mechanisms2025projective} showed that Equation~\ref{eq:poe_score} is only
well-defined --- in the sense that its samples remain on the data manifold --- when the
conditional score fields act on approximately orthogonal directions in the data representation.
Formally, let $A$ be a linear transformation of the data. The \emph{Factorised Conditionals}
condition requires that in the transformed space, the score contributions from $c_1$ and $c_2$
act on separable, approximately non-overlapping subspaces:
\begin{equation}
    \nabla_x \log p(x \mid c_i) - \nabla_x \log p(x)
    \;\perp_A\;
    \nabla_x \log p(x \mid c_j) - \nabla_x \log p(x).
    \label{eq:factorised_conditionals}
\end{equation}
When this condition holds, the composed score is a valid score function and denoising under
it converges to a coherent joint sample. When it fails --- either because the concepts
compete for the same feature dimensions or because the joint configuration $p(x \mid c_1,
c_2)$ is absent from the training distribution --- the composed score can be incoherent or
degenerate.

A critical consequence is that PoE does \emph{not} generalise out of distribution. When no
training images jointly exhibit both $c_1$ and $c_2$, Equation~\ref{eq:poe_dist} does not
define a new distribution; it inherits the support of the component models, and the composed
samples are constrained to lie near the modes of whichever single-concept prior dominates.
This is the theoretical basis for the composability gap studied in Phase~1: composition
succeeds when the joint configuration is represented in the training data (co-occurring
concepts, Group~1) or when the concepts act on orthogonal feature dimensions (disentangled
concepts, Group~2), and fails when neither condition holds (Groups~3 and~4).

\subsection{From Score Composition to Representation}
\label{subsec:bg_representation}

Du et al.~\citep{du2024compositional} showed that a \emph{single} pretrained model can be
composed at inference time by running multiple conditional passes and combining the resulting
score fields, without requiring a separately trained expert for each concept. Their analysis
establishes that inference-time choices --- guidance weighting, sampling schedule, and
step-size --- affect the stability and quality of composition, and they propose dynamic
guidance schedules to mitigate mode collapse under strong guidance. Crucially, however,
these inference-time interventions do not alter the underlying conditional independence
assumption (Equation~\ref{eq:poe_dist}): they improve how well the operator is applied
given the existing score fields, but cannot compensate for score fields that do not
satisfy the Factorised Conditionals condition.

This observation motivates a representational account of compositional failure. If the
learned features that the score function acts on are entangled --- if the representation of
one concept overlaps geometrically with the representation of another --- then no choice of
composition operator or guidance schedule can cleanly separate them at the score level. The
gap is not a \emph{rule problem} (which operator to use) but a \emph{representation problem}
(whether the score function has learned a factorised encoding of the concepts). Phase~2 tests
this directly by holding the composition rule fixed and varying only the learned
representation, asking whether progressively more factorised latent codes reduce the same
compositional failures observed in Phase~1.
